% GENERATED FILE. Edit canonical lessons, then run npm run generate:books.
% canonical-source-hash: fnv1a64:b65559f5236463b0
% canonical-lessons: HI-C54-cow, HI-C54-horse, HI-C54-fish, HI-C54-bird, HI-C54-goat

\chapter{Animals}
\label{ch:hi-animals}
\hlchaptermodality{54}

\begin{chapteropening}
\textbf{By the end of this chapter:} \emph{I can name five common animals when I say what something is called.}

Complete HI-C54-goat and say all five words in order.
\end{chapteropening}

\section[gāy]{\hi{गाय} (gāy) — a cow}

\label{lesson:HI-C54-cow}



\begin{quote}
Before the new run of words: what did the last two words of the chapter before mean?
\end{quote}

\subsection*{You'll want to know: \hi{गाय}}
\textbf{\hi{गाय}} (\emph{gāy}) — \textquotedblleft{}a cow\textquotedblright{}.

From Sanskrit \hi{गो} (\emph{go}), 'cow' — and this one goes all the way back. The reconstructed ancestor is \emph{g\textsuperscript{w}\'{\={o}}ws}, which gives Latin \emph{bōs}, Greek \emph{boûs}, and English \emph{cow}.

So \hi{गाय} and \emph{cow} are the same word, separated by five thousand years and most of a continent. It is one of the cleanest cousins in this whole book: no borrowing, no detour, just two ends of one inheritance.

\begin{grammarlens}[title={What you've built}]
The first of five animals, and a word English owns too.
\end{grammarlens}

\subsection*{Your turn}
Say these aloud:

\begin{itemize}
\raggedright
  \item \emph{gāy}
  \item it once more, slowly
  \item \emph{gāy}, and hear the \emph{c-} of \emph{cow} inside it
\end{itemize}

\begin{tcolorbox}[breakable,colback=teal!4,colframe=teal!35!black,title={Before you move on}]
What does \hi{गाय} mean? (\textquotedblleft{}a cow\textquotedblright{}.) Say it once more.
\end{tcolorbox}
\section[ghoṛā]{\hi{घोड़ा} (ghoṛā) — a horse}

\label{lesson:HI-C54-horse}



\begin{quote}
What did the word before this one mean?
\end{quote}

\subsection*{You'll want to know: \hi{घोड़ा}}
\textbf{\hi{घोड़ा}} (\emph{ghoṛā}) — \textquotedblleft{}a horse\textquotedblright{}.

From Sanskrit \hi{घोटक} (\emph{ghoṭaka}), 'horse' — a word whose own origin is uncertain and may not be Indo-European at all.

What makes it worth stopping on is what did \emph{not} happen. Sanskrit had \hi{अश्व} (\emph{aśva}), the inherited horse-word, a straight cousin of Latin \emph{equus} and Greek \emph{híppos}. It lost. Everyday Hindi rides \hi{घोड़ा}, and the ancient word survives only in names and formal compounds.

\begin{grammarlens}[title={What you've built}]
The second of five, and an inherited word that was pushed out.
\end{grammarlens}

\subsection*{Your turn}
Say these aloud:

\begin{itemize}
\raggedright
  \item \emph{ghoṛā}
  \item it once more, slowly
  \item \emph{ghoṛā}, then \emph{gāy} — the one that was replaced, the one that was not
\end{itemize}

\begin{tcolorbox}[breakable,colback=teal!4,colframe=teal!35!black,title={Before you move on}]
What does \hi{घोड़ा} mean? (\textquotedblleft{}a horse\textquotedblright{}.) Say it once more.
\end{tcolorbox}
\section[machhlī]{\hi{मछली} (machhlī) — a fish}

\label{lesson:HI-C54-fish}



\begin{quote}
Name the 2 words of this chapter so far.
\end{quote}

\subsection*{You'll want to know: \hi{मछली}}
\textbf{\hi{मछली}} (\emph{machhlī}) — \textquotedblleft{}a fish\textquotedblright{}.

From Sanskrit \hi{मत्स्य} (\emph{matsya}), 'fish', through a Prakrit form in which the \emph{-tsy-} cluster fused into the \emph{-chh-} you now say.

The \emph{-ī} on the end is the same feminine ending you have met on \hi{कुर्सी} and \hi{बिल्ली}, and \hi{मछली} is feminine like them. The ending is doing grammatical work here, not naming anything.

\begin{grammarlens}[title={What you've built}]
The third of five, and another cluster that gave way.
\end{grammarlens}

\subsection*{Your turn}
Say these aloud:

\begin{itemize}
\raggedright
  \item \emph{machhlī}
  \item it once more, slowly
  \item \emph{machhlī}, and remember it is feminine
\end{itemize}

\begin{tcolorbox}[breakable,colback=teal!4,colframe=teal!35!black,title={Before you move on}]
What does \hi{मछली} mean? (\textquotedblleft{}a fish\textquotedblright{}.) Say it once more.
\end{tcolorbox}
\section[chiṛiyā]{\hi{चिड़िया} (chiṛiyā) — a small bird}

\label{lesson:HI-C54-bird}



\begin{quote}
Name the 3 words of this chapter so far.
\end{quote}

\subsection*{You'll want to know: \hi{चिड़िया}}
\textbf{\hi{चिड़िया}} (\emph{chiṛiyā}) — \textquotedblleft{}a small bird\textquotedblright{}.

The origin of \hi{चिड़िया} is not settled. It is usually grouped with a family of words for small birds that look as though they were built from the noise a bird makes, but that is a description of the family, not a documented derivation.

It is worth saying plainly that 'sounds like a bird' is a guess and not evidence. The honest position is that this everyday word has no agreed ancestor, and that is the fifth such word you have met.

\begin{grammarlens}[title={What you've built}]
The fourth of five, and a guess that this book will not make.
\end{grammarlens}

\subsection*{Your turn}
Say these aloud:

\begin{itemize}
\raggedright
  \item \emph{chiṛiyā}
  \item it once more, slowly
  \item \emph{chiṛiyā}, small and feminine
\end{itemize}

\begin{tcolorbox}[breakable,colback=teal!4,colframe=teal!35!black,title={Before you move on}]
What does \hi{चिड़िया} mean? (\textquotedblleft{}a small bird\textquotedblright{}.) Say it once more.
\end{tcolorbox}
\section[bakrī]{\hi{बकरी} (bakrī) — a goat}

\label{lesson:HI-C54-goat}



\begin{quote}
Name the 4 words of this chapter so far.
\end{quote}

\subsection*{You'll want to know: \hi{बकरी}}
\textbf{\hi{बकरी}} (\emph{bakrī}) — \textquotedblleft{}a goat\textquotedblright{}.

From Sanskrit \hi{बर्कर} (\emph{barkara}), 'a kid, a young goat'. The masculine \hi{बकरा} and feminine \hi{बकरी} split the animal by sex, as Hindi does with \hi{बेटा} and \hi{बेटी}.

Notice what the pair does that English does not: \emph{goat} covers both, and English has to add a word to say which. Hindi has already said it in the ending.

\begin{grammarlens}[title={What you've built}]
The fifth of five. That is the yard.
\end{grammarlens}

\subsection*{Your turn}
Say these aloud:

\begin{itemize}
\raggedright
  \item \emph{bakrī}
  \item it once more, slowly
  \item \emph{bakrī}, \emph{gāy}, \emph{ghoṛā}, \emph{machhlī}, \emph{chiṛiyā} — name them round
\end{itemize}

\begin{tcolorbox}[breakable,colback=teal!4,colframe=teal!35!black,title={Before you move on}]
What does \hi{बकरी} mean? (\textquotedblleft{}a goat\textquotedblright{}.) Say it once more.
\end{tcolorbox}
